\documentclass[11pt]{article}
% Shared preamble for the printable mini lessons (1.4 onward).
%   \documentclass[11pt]{article}
%   % Shared preamble for the printable mini lessons (1.4 onward).
%   \documentclass[11pt]{article}
%   % Shared preamble for the printable mini lessons (1.4 onward).
%   \documentclass[11pt]{article}
%   \input{../preamble}
% Build the key with:  pdflatex -jobname=<name>_KEY "\def\iskey{}\input{<name>}"
% Figures shared with the web page live in figures/Unit_X/<name>.tex and are
% pulled in with \fig{Unit_X/<name>}; the standalone package drops their own
% preamble, so their macros come from figures/preamble.tex, loaded here.
\input{../../figures/preamble}
\usepackage[margin=0.75in]{geometry}
\usepackage{array}
\usepackage{enumitem}
\usepackage{amsmath}
\usepackage[helvet]{sfmath}  % math in the same sans font as the text
\usepackage{standalone}
\pagestyle{empty}
\setlength{\parindent}{0pt}
\setlist[enumerate,1]{leftmargin=1.5em,itemsep=1em}
\setlist[enumerate,2]{label=(\alph*),leftmargin=2em,itemsep=2.2em}

\newif\ifkey
\ifdefined\iskey \keytrue \fi
\colorlet{keyred}{red!75!black}
\newcommand{\ans}[1]{\ifkey\leavevmode{\color{keyred}#1}\fi}
% vertical work space on the student copy, less on the key
\newcommand{\work}[1]{\ifkey\vspace{0.6em}\else\vspace{#1}\fi}
\newcommand{\fig}[1]{\input{../../figures/#1}}
% Title line: \lessontitle{Quartiles, IQR, and Box Plots}
\newcommand{\lessontitle}[1]{{\Large\bfseries #1\ifkey\ \textcolor{keyred}{KEY}\fi}\hfill Name \underline{\hspace{2in}}\par\vspace{0.6em}}
\newcommand{\partner}{\par\vspace{0.4em}Work with a partner. Compare answers when you finish each problem.\par}

% Build the key with:  pdflatex -jobname=<name>_KEY "\def\iskey{}\input{<name>}"
% Figures shared with the web page live in figures/Unit_X/<name>.tex and are
% pulled in with \fig{Unit_X/<name>}; the standalone package drops their own
% preamble, so their macros come from figures/preamble.tex, loaded here.
% Shared preamble for web figures.  Each figure is a standalone .tex that
% \input's this file; scripts/make-figures.sh turns it into an SVG.
\usepackage[T1]{fontenc}
\usepackage{helvet}
\renewcommand{\familydefault}{\sfdefault}
\usepackage{tikz}
\usetikzlibrary{decorations.pathreplacing}
% axis numbers in the text font, not math italic
\pgfkeys{/pgf/number format/.cd,assume math mode=true,1000 sep={}}

% dot plot: \dotplot{min}{max}{value/count,...}{axis label}
\newcommand{\dotplot}[4]{%
\begin{tikzpicture}[x=0.9cm,y=0.33cm]
  \draw[thick] (#1-0.5,0) -- (#2+0.5,0);
  \foreach \t in {#1,...,#2} {\draw (\t,0.25) -- (\t,-0.25) node[below] {\t};}
  \foreach \v/\c in {#3} {\ifnum\c>0 \foreach \k in {1,...,\c} {\fill (\v,\k*1.15) circle (0.16cm);}\fi}
  \node at ({(#1+#2)/2},-2.3) {#4};
\end{tikzpicture}}

% shape sketch histogram (12 bins): \shapehist{heights}{label}
\newcommand{\shapehist}[2]{%
  \foreach \h [count=\i from 0] in {#1} \draw[fill=gray!30] (\i,0) rectangle (\i+1,\h);
  \draw (-0.3,0) -- (12.3,0);
  \node[font=\bfseries] at (6,-2.6) {#2};}
\newcommand{\meanmark}[1]{\fill (#1,0) -- ++(-0.3,-0.75) -- ++(0.6,0) -- cycle;}
\newcommand{\medmark}[2]{\draw[very thick,dashed] (#1,0) -- (#1,#2);}

% number line with minor ticks: \nline{min}{max}{label step}{axis label}
\newcommand{\nline}[4]{%
  \draw[thick] (#1-0.5,0) -- (#2+0.5,0);
  \foreach \t in {#1,...,#2} \draw (\t,0.12) -- (\t,-0.12);
  \pgfmathtruncatemacro{\nxt}{#1+#3}
  \foreach \t in {#1,\nxt,...,#2} \draw (\t,0.2) -- (\t,-0.2) node[below] {\t};
  \node at ({(#1+#2)/2},-0.95) {#4};}
% box plot above the line: \bplot{min}{Q1}{med}{Q3}{max}
\newcommand{\bplot}[5]{%
  \draw[thick] (#1,1.0) -- (#2,1.0);
  \draw[thick] (#4,1.0) -- (#5,1.0);
  \draw[thick] (#1,0.75) -- (#1,1.25);
  \draw[thick] (#5,0.75) -- (#5,1.25);
  \draw[thick,fill=gray!20] (#2,0.55) rectangle (#4,1.45);
  \draw[very thick] (#3,0.55) -- (#3,1.45);}

% ---- added for 1.4 onward ----
% axis on y=0: \axis{min}{max}{label step}{tick step}{axis label}
\newcommand{\axis}[5]{%
  \draw[thick] (#1,0) -- (#2,0);
  \pgfmathsetmacro{\axT}{#1+#4}\foreach \t in {#1,\axT,...,#2} \draw (\t,0.1) -- (\t,-0.1);
  \pgfmathsetmacro{\axL}{#1+#3}\foreach \t in {#1,\axL,...,#2}
    \draw (\t,0.18) -- (\t,-0.18) node[below] {\pgfmathprintnumber{\t}};
  \node at ({(#1+#2)/2},-0.95) {#5};}
% box plot centred at height y: \bplotat{y}{min}{Q1}{med}{Q3}{max}
\newcommand{\bplotat}[6]{%
  \draw[thick] (#2,#1) -- (#3,#1);
  \draw[thick] (#5,#1) -- (#6,#1);
  \draw[thick] (#2,#1-0.22) -- (#2,#1+0.22);
  \draw[thick] (#6,#1-0.22) -- (#6,#1+0.22);
  \draw[thick,fill=gray!20] (#3,#1-0.42) rectangle (#5,#1+0.42);
  \draw[very thick] (#4,#1-0.42) -- (#4,#1+0.42);}
% an outlier dot: \outlier{x}{y}
\newcommand{\outlier}[2]{\fill (#1,#2) circle (2.4pt);}

\usepackage[margin=0.75in]{geometry}
\usepackage{array}
\usepackage{enumitem}
\usepackage{amsmath}
\usepackage[helvet]{sfmath}  % math in the same sans font as the text
\usepackage{standalone}
\pagestyle{empty}
\setlength{\parindent}{0pt}
\setlist[enumerate,1]{leftmargin=1.5em,itemsep=1em}
\setlist[enumerate,2]{label=(\alph*),leftmargin=2em,itemsep=2.2em}

\newif\ifkey
\ifdefined\iskey \keytrue \fi
\colorlet{keyred}{red!75!black}
\newcommand{\ans}[1]{\ifkey\leavevmode{\color{keyred}#1}\fi}
% vertical work space on the student copy, less on the key
\newcommand{\work}[1]{\ifkey\vspace{0.6em}\else\vspace{#1}\fi}
\newcommand{\fig}[1]{\input{../../figures/#1}}
% Title line: \lessontitle{Quartiles, IQR, and Box Plots}
\newcommand{\lessontitle}[1]{{\Large\bfseries #1\ifkey\ \textcolor{keyred}{KEY}\fi}\hfill Name \underline{\hspace{2in}}\par\vspace{0.6em}}
\newcommand{\partner}{\par\vspace{0.4em}Work with a partner. Compare answers when you finish each problem.\par}

% Build the key with:  pdflatex -jobname=<name>_KEY "\def\iskey{}\input{<name>}"
% Figures shared with the web page live in figures/Unit_X/<name>.tex and are
% pulled in with \fig{Unit_X/<name>}; the standalone package drops their own
% preamble, so their macros come from figures/preamble.tex, loaded here.
% Shared preamble for web figures.  Each figure is a standalone .tex that
% \input's this file; scripts/make-figures.sh turns it into an SVG.
\usepackage[T1]{fontenc}
\usepackage{helvet}
\renewcommand{\familydefault}{\sfdefault}
\usepackage{tikz}
\usetikzlibrary{decorations.pathreplacing}
% axis numbers in the text font, not math italic
\pgfkeys{/pgf/number format/.cd,assume math mode=true,1000 sep={}}

% dot plot: \dotplot{min}{max}{value/count,...}{axis label}
\newcommand{\dotplot}[4]{%
\begin{tikzpicture}[x=0.9cm,y=0.33cm]
  \draw[thick] (#1-0.5,0) -- (#2+0.5,0);
  \foreach \t in {#1,...,#2} {\draw (\t,0.25) -- (\t,-0.25) node[below] {\t};}
  \foreach \v/\c in {#3} {\ifnum\c>0 \foreach \k in {1,...,\c} {\fill (\v,\k*1.15) circle (0.16cm);}\fi}
  \node at ({(#1+#2)/2},-2.3) {#4};
\end{tikzpicture}}

% shape sketch histogram (12 bins): \shapehist{heights}{label}
\newcommand{\shapehist}[2]{%
  \foreach \h [count=\i from 0] in {#1} \draw[fill=gray!30] (\i,0) rectangle (\i+1,\h);
  \draw (-0.3,0) -- (12.3,0);
  \node[font=\bfseries] at (6,-2.6) {#2};}
\newcommand{\meanmark}[1]{\fill (#1,0) -- ++(-0.3,-0.75) -- ++(0.6,0) -- cycle;}
\newcommand{\medmark}[2]{\draw[very thick,dashed] (#1,0) -- (#1,#2);}

% number line with minor ticks: \nline{min}{max}{label step}{axis label}
\newcommand{\nline}[4]{%
  \draw[thick] (#1-0.5,0) -- (#2+0.5,0);
  \foreach \t in {#1,...,#2} \draw (\t,0.12) -- (\t,-0.12);
  \pgfmathtruncatemacro{\nxt}{#1+#3}
  \foreach \t in {#1,\nxt,...,#2} \draw (\t,0.2) -- (\t,-0.2) node[below] {\t};
  \node at ({(#1+#2)/2},-0.95) {#4};}
% box plot above the line: \bplot{min}{Q1}{med}{Q3}{max}
\newcommand{\bplot}[5]{%
  \draw[thick] (#1,1.0) -- (#2,1.0);
  \draw[thick] (#4,1.0) -- (#5,1.0);
  \draw[thick] (#1,0.75) -- (#1,1.25);
  \draw[thick] (#5,0.75) -- (#5,1.25);
  \draw[thick,fill=gray!20] (#2,0.55) rectangle (#4,1.45);
  \draw[very thick] (#3,0.55) -- (#3,1.45);}

% ---- added for 1.4 onward ----
% axis on y=0: \axis{min}{max}{label step}{tick step}{axis label}
\newcommand{\axis}[5]{%
  \draw[thick] (#1,0) -- (#2,0);
  \pgfmathsetmacro{\axT}{#1+#4}\foreach \t in {#1,\axT,...,#2} \draw (\t,0.1) -- (\t,-0.1);
  \pgfmathsetmacro{\axL}{#1+#3}\foreach \t in {#1,\axL,...,#2}
    \draw (\t,0.18) -- (\t,-0.18) node[below] {\pgfmathprintnumber{\t}};
  \node at ({(#1+#2)/2},-0.95) {#5};}
% box plot centred at height y: \bplotat{y}{min}{Q1}{med}{Q3}{max}
\newcommand{\bplotat}[6]{%
  \draw[thick] (#2,#1) -- (#3,#1);
  \draw[thick] (#5,#1) -- (#6,#1);
  \draw[thick] (#2,#1-0.22) -- (#2,#1+0.22);
  \draw[thick] (#6,#1-0.22) -- (#6,#1+0.22);
  \draw[thick,fill=gray!20] (#3,#1-0.42) rectangle (#5,#1+0.42);
  \draw[very thick] (#4,#1-0.42) -- (#4,#1+0.42);}
% an outlier dot: \outlier{x}{y}
\newcommand{\outlier}[2]{\fill (#1,#2) circle (2.4pt);}

\usepackage[margin=0.75in]{geometry}
\usepackage{array}
\usepackage{enumitem}
\usepackage{amsmath}
\usepackage[helvet]{sfmath}  % math in the same sans font as the text
\usepackage{standalone}
\pagestyle{empty}
\setlength{\parindent}{0pt}
\setlist[enumerate,1]{leftmargin=1.5em,itemsep=1em}
\setlist[enumerate,2]{label=(\alph*),leftmargin=2em,itemsep=2.2em}

\newif\ifkey
\ifdefined\iskey \keytrue \fi
\colorlet{keyred}{red!75!black}
\newcommand{\ans}[1]{\ifkey\leavevmode{\color{keyred}#1}\fi}
% vertical work space on the student copy, less on the key
\newcommand{\work}[1]{\ifkey\vspace{0.6em}\else\vspace{#1}\fi}
\newcommand{\fig}[1]{\input{../../figures/#1}}
% Title line: \lessontitle{Quartiles, IQR, and Box Plots}
\newcommand{\lessontitle}[1]{{\Large\bfseries #1\ifkey\ \textcolor{keyred}{KEY}\fi}\hfill Name \underline{\hspace{2in}}\par\vspace{0.6em}}
\newcommand{\partner}{\par\vspace{0.4em}Work with a partner. Compare answers when you finish each problem.\par}

% Unit summary: the unit's tools in tables, to print and keep. Same tables as
% chapters/Unit_4/Summary.qmd; edit both together. Nothing to answer, so the
% _KEY copy is the same sheet.
\newcommand{\sumhead}[1]{\par\vspace{0.7em}{\large\bfseries #1}\par\vspace{0.3em}}
\newcolumntype{L}[1]{>{\raggedright\arraybackslash}p{#1}}
\renewcommand{\arraystretch}{1.25}
% room above and below a full-size fraction so it clears the rules
\newcommand{\fracstrut}{\rule[-0.95em]{0pt}{2.3em}}

% Same title on the student and _KEY copies: a summary has no answers, so the
% key is the same sheet and must not grow a line for the word KEY.
\newlength{\sumtw}
\newcommand{\sumtitle}[1]{\settowidth{\sumtw}{\Large\bfseries #1}%
  \ifdim\sumtw>4.3in \hfill Name \underline{\hspace{2in}}\par\vspace{0.5em}{\Large\bfseries #1}\par\vspace{0.6em}%
  \else {\Large\bfseries #1}\hfill Name \underline{\hspace{2in}}\par\vspace{0.6em}\fi}

\begin{document}

\sumtitle{Unit 4 Summary: Scatterplots}

\sumhead{Describing a Scatterplot}
A \textbf{scatterplot} shows two numbers for each person or thing. The first number ($x$) goes across. The second number ($y$) goes up.

\vspace{0.4em}
\begin{tabular}{|L{1.2in}|L{2.75in}|L{2.5in}|}\hline
\textbf{Feature} & \textbf{Ask} & \textbf{Example: the study data}\\ \hline
Direction & As $x$ goes up, does $y$ go up (\textbf{positive}) or down (\textbf{negative})? No pattern at all is \textbf{no association}. & \textbf{Positive.} The students who studied more mostly scored higher.\\ \hline
Form & Do the dots follow a line (\textbf{linear}) or a curve? & \textbf{Linear.} The dots follow a line.\\ \hline
Strength & Are the dots close to a line (\textbf{strong}), a little spread out (\textbf{moderate}), or very spread out (\textbf{weak})? & \textbf{Strong.} The dots are close to the line.\\ \hline
Unusual & Any outliers or clusters? & None.\\ \hline
\end{tabular}

\sumhead{The Line of Best Fit}
A \textbf{line of best fit} follows the trend of the dots. It goes through the middle of the dots. About the same number of dots are above it as below it.

\vspace{0.4em}
\begin{tabular}{|L{1.2in}|L{2.75in}|L{2.5in}|}\hline
\textbf{Step} & \textbf{How it works} & \textbf{Example: the study data}\\ \hline
1. Pick two points & Pick two points \textbf{on the line}. They do not have to be dots. & $(2, 64)$ and $(6, 84)$\\ \hline
2. Find the slope~$m$ & Subtract the $y$ values. Subtract the $x$ values in the same order. Divide. & \fracstrut$m = \dfrac{84 - 64}{6 - 2} = \dfrac{20}{4} = 5$\\ \hline
3. Find the y-intercept~$b$ & The equation is $\hat{y} = mx + b$. Put in $m$ and one point. Solve for $b$. & $64 = 5(2) + b$, \; so $b = 54$\\ \hline
4. Write the equation & Put $m$ and $b$ into $\hat{y} = mx + b$. $\hat{y}$ is read ``y-hat.'' The hat means \textbf{predicted}. & $\hat{y} = 5x + 54$\\ \hline
5. Predict & Put the $x$ value into the equation. The line predicts $y$. It does not give the exact $y$. & A student studies 5 hours. $\hat{y} = 5(5) + 54 = 79$. The predicted score is \textbf{79}.\\ \hline
Extrapolation & Predicting far outside the data. It is risky. & 12 hours: $\hat{y} = 5(12) + 54 = 114$. A test score cannot be more than 100. The data only goes from 1 to 8 hours.\\ \hline
\end{tabular}

\sumhead{Slope and Intercept in Context}
\begin{tabular}{|L{0.8in}|L{2.1in}|L{1.7in}|L{1.7in}|}\hline
 & \textbf{Sentence frame} & \textbf{A plant: $\hat{y} = 2.5x + 4$} & \textbf{A candle: $\hat{y} = -2x + 24$}\\ \hline
Slope & For each additional [$x$ unit], the predicted [$y$] increases (or decreases) by [slope and units]. & For each additional \textbf{week}, the predicted \textbf{height} increases by \textbf{2.5 cm}. & For each additional \textbf{hour}, the predicted \textbf{height} decreases by \textbf{2 cm}.\\ \hline
y-intercept & When [$x$] is 0, the predicted [$y$] is [intercept and units]. & When $x = 0$ (the day it was planted), the predicted height is \textbf{4 cm}. & When $x = 0$ (a new candle), the predicted height is \textbf{24 cm}.\\ \hline
\end{tabular}

\vspace{0.3em}
\textbf{Some intercepts have no meaning:} if $x$ is a person's height, then $x = 0$ is not possible.

\newpage
\sumhead{Correlation and Causation}
\begin{tabular}{|L{1.2in}|L{2.75in}|L{2.5in}|}\hline
\textbf{Idea} & \textbf{What it means} & \textbf{Example}\\ \hline
Correlation coefficient $r$ & It measures how strong a \textbf{linear} pattern is. $r$ is always from $-1$ to $1$. & The candles: $r \approx -0.98$. The dots are close to a line that goes down.\\ \hline
Sign of $r$ & The sign shows the direction. $r > 0$ is positive. $r < 0$ is negative. & $r \approx 0.95$ is positive. $r \approx -0.80$ is negative.\\ \hline
Size of $r$ & Close to $1$ or $-1$ is strong. Close to $0$ is weak. & $r = -0.8$ is stronger than $r = 0.6$. It is closer to $-1$ than $0.6$ is to $1$. The minus sign only tells the direction.\\ \hline
Correlation is not causation & A correlation shows that two things change together. It does \textbf{not} prove that one causes the other. & Ice cream sales and sunburns have a strong positive correlation. Ice cream does not cause sunburns.\\ \hline
Lurking variable & A hidden cause of both things. & Hot, sunny weather causes \textbf{both}.\\ \hline
\end{tabular}

\sumhead{Residuals}
\begin{tabular}{|L{1.2in}|L{2.75in}|L{2.5in}|}\hline
\textbf{Idea} & \textbf{What it means} & \textbf{Example: the study line $\hat{y} = 5x + 54$}\\ \hline
Residual & residual $=$ actual $-$ predicted $= y - \hat{y}$ & A student studied 4 hours and scored 78. $\hat{y} = 5(4) + 54 = 74$. Residual $= 78 - 74 = 4$.\\ \hline
Positive residual & The point is \textbf{above} the line. & Residual 4: the student scored 4 points \textbf{more} than predicted.\\ \hline
Negative residual & The point is \textbf{below} the line. & Another student studied 4 hours and scored 72. Residual $= 72 - 74 = -2$. The student scored 2 points \textbf{less} than predicted.\\ \hline
\end{tabular}

\sumhead{Least Squares and Residual Plots}
\begin{tabular}{|L{1.2in}|L{2.75in}|L{2.5in}|}\hline
\textbf{Tool} & \textbf{How it works} & \textbf{What it means}\\ \hline
Least-squares regression line & Square every residual and add them up: $\Sigma(y - \hat{y})^2$. $\Sigma$ means add up. & This line makes the sum as \textbf{small as possible}. Calculators and computers find this line for you.\\ \hline
Residual plot & Put $x$ across the bottom and the residuals up the side. & Random scatter, no pattern: a line is a \textbf{good} model. A curved pattern: a line is \textbf{not} a good model. The data is curved.\\ \hline
\end{tabular}

\end{document}
